\clearpage
\clearpage
\thispagestyle{empty}

\vspace*{\fill}

\begin{center}
\begin{minipage}{0.85\textwidth}

\centering

\vspace{1em}

\raggedright
I want to express my gratitude to my family for their never ending support from back home. Their motivation and belief in me has pushed me to make this paper the best it can be. I could not have pushed through this degree without them by my side.
\\ \vspace{1em}
I would also like to thank Professor Vincent Van Eylen for allowing me to take on this project. The guidance from him along side his PhD student Marina \textbf{(lastname)} and PDRA Kendall Sullivan has been invaluable in coming to the results in this paper. \textbf{(Do I need to provide proper titles for everyone?)}. My appreciation for astrophysics research has only grown with this project and I appreciate their help along the way.

\end{minipage}
\end{center}

\vspace*{\fill}

\clearpage

\newpage
\addcontentsline{toc}{section}{Abstract}
\begin{abstract}
    The use of convolutional neural networks (CNN) to analyse patterns in exoplanet transit signals has been  widely applied for identifying false positive signals. Previous studies commonly focus on substantial preprocessing of the data to improve classification from the model. This paper focuses on improving a CNN model to classify Kepler light curves with limited preprocessing for classification. This included avoiding phase-folding, binning, fitting and other typical tasks used for preparing light curve data. Using lightly processed data, one binary CNN was developed to classify Kepler transiting exoplanets against eclipsing binaries, a common false positive in light curve analysis. A multiclass CNN was also developed to classify light curves from Kepler as transits, shallow or deep eclipsing binaries where a boundary between the eclipsing binary groups was developed for the operations in this work. The binary model evaluation achieved a 92\% accuracy score with 96\% recall indicating the model correctly identified almost all transiting exoplanets. The multiclass model had an accuracy of 78\% and macro averaged F1 score of 79\%. The main source of model confusion comes from the classification difficulty of shallow eclipsing binaries which can mimic exoplanet transit signals. Two random forest models, one binary and one multiclass, were developed as a baseline to compare the performance of the neural networks. The comparison showed the importance of the statistical features of the light curves are as the CNN models achieved similar performance to the random forest models. 
\end{abstract}

\newpage
\tableofcontents
\newpage
\pagenumbering{arabic}
\section{Introduction}

Prior to 1992, the existence of planets beyond the Solar System was largely speculative. This changed with the first observational evidence of extrasolar objects in the early 1990s \citep{First_Exoplanet, First_Exoplanet_RV}. These confirmations transformed the field from theoretical speculation into an observational science. Since then, the rate of discovery has increased rapidly, and by 2008 over 263 exoplanet systems had been confirmed \citep{Exoplanet_numbers_2008}. As of August 2026, over 6000 exoplanets have been confirmed, and thousands more are candidates awaiting further observation \citep{christiansen2025nasaexoplanetarchiveexoplanet}.\footnote{Confirmed planet counts obtained from the NASA Exoplanet Archive Confirmed Planets Table, accessed March 2026: \url{https://exoplanetarchive.ipac.caltech.edu/}}

This rapid growth in detections has been driven by advances in observational techniques and instrumentation. 51 Pegasi b was the first exoplanet discovered orbiting a Sun-like star. It was detected using radial velocity (RV) measurements of its host star \citep{First_Exoplanet_RV}. Although radial velocity measurements remain widely used, transit photometry has become one of the most successful detection methods \citep{How_Many_Transits}. With the development of large-format CCDs, the transit method became the most practical way for large scale exoplanets surveys as wide-field observations allow large numbers of stars to be monitored simultaneously \citep{2005MNRAS.360..703H}.

When analysing transit data, a major source of false positive transit-like signals is eclipsing binary systems, where two stars orbit one another and periodically pass in front of each other from the observer's perspective \citep{2018ApJS..235...38T, Morton_2016}. These systems can produce repeated decreases in observed flux that resemble planetary transits, making them an important contaminant when identifying exoplanet candidates.

This paper trains machine learning models to classify light curves collected from NASA's \textit{Kepler} mission as either transiting exoplanets or eclipsing binaries. Since eclipsing binaries and other astrophysical systems can produce transit-like false positive signals, an automated classifier could reduce the need for manual vetting. Previous studies have explored machine learning approaches for false positive identification, often using heavily processed or phase-folded light curves. The goal of this work is to test whether a model can achieve useful classification performance using relatively limited preprocessing of the input data.

